% =====================================================================
%  PeerJ submission — SUPPLEMENTARY MATERIAL
%  EEG–Hemodynamic Coupling (Topic 9, null result)
%  Companion to NCCI_peerj.tex. Compiles standalone with wlpeerj.cls.
%  Figures live one directory up (manuscript/).
% =====================================================================
\documentclass[fleqn,10pt]{wlpeerj}

\usepackage{float}
\usepackage{array}
\graphicspath{{../}}

\usepackage[breaklinks=true]{hyperref}
\definecolor{OldCyan}{RGB}{62,127,137}
\hypersetup{colorlinks,citecolor=OldCyan,filecolor=OldCyan,
            linkcolor=OldCyan,urlcolor=OldCyan}

\setcounter{secnumdepth}{0}

% Number supplementary items as S1, S2, ...
\renewcommand{\thesection}{S\arabic{section}}
\renewcommand{\thefigure}{S\arabic{figure}}
\renewcommand{\thetable}{S\arabic{table}}
\renewcommand{\theequation}{S\arabic{equation}}

\title{Supplementary Material\\[4pt]
\large EEG--Hemodynamic Coupling During General Anaesthesia:
A Large-Scale Exploratory Analysis in 2,070 Surgical Cases
with Null Acute Event-Response Findings}

\author[1]{Ge Gao}
\affil[1]{Department of Anesthesiology, The First Affiliated Hospital,
Zhejiang University School of Medicine, Hangzhou, China}

\keywords{supplementary material}

\begin{abstract}
This document contains Supplementary Text S1--S6 and Supplementary Figures S1--S3
referenced in the main manuscript. Section, figure, table and equation numbers are
prefixed with ``S''.
\end{abstract}

\begin{document}

\flushbottom
\maketitle
\thispagestyle{empty}

% ── S1 ───────────────────────────────────────────────────────────────────────
\section{S1. NCCI and WCS Construction}

\textbf{NCCI (global PCA).} Input per case: [GC$_\text{e2h}$ F (winsorised),
WC$_\text{overall}$, Pearson $r$], $z$-score standardised. A single PCA was fit across all
2,070 cases: PC1 37.0\%, PC2 33.9\%, PC3 29.1\% of variance. PC1 was MinMax-scaled to [0,1].
Sign convention: PC1 inverted if corr(PC1,~WC)$<$0.

\textbf{WCS (per-case local PCA).} Per case: per-window metrics [GC$_\text{e2h}$ F, WC,
Pearson $r$], $z$-score standardised on case-level mean/SD. A single-component PCA was fit
per case. Sign: same rule (inverted if corr(PC1,~WC)$<$0). MinMax within case. PC1 variance
$\approx$37--42\% (variable across cases).

\textbf{GEE justification.} Modelling $\Delta$WCS (the within-case relative change) means the
outcome is already patient-centred. The GEE exchangeable structure handles the residual
clustering of events within the same patient. WCS does not need to be on a common scale across
patients for this within-case difference to be meaningful; the sign-alignment step ensures
that $\Delta$WCS $>$0 always denotes the same physiological direction (higher coupling) within
any patient.

% ── S2 ───────────────────────────────────────────────────────────────────────
\section{S2. GEE Implementation}

\texttt{statsmodels.formula.api.gee}; formula: \texttt{delta\_wcs $\sim$ 1}; groups: patient
ID; family: Gaussian; cov\_struct: Exchangeable (primary) / Autoregressive (sensitivity);
robust sandwich variance. $n_\text{obs} = 73{,}001$; groups: 2,031.

% ── S3 ───────────────────────────────────────────────────────────────────────
\section{S3. Granger Causality Implementation}

\texttt{statsmodels.tsa.stattools.grangercausalitytests}, maxlag $= 5$. ADF stationarity test
($\alpha = 0.05$); first-difference if non-stationary. Lag selection: minimum SSR F-test
$p$-value. Windows with $<$20 valid points were excluded. Winsorisation at the 99th percentile
($\hat{c} = 4.26$) was applied for PCA input; raw values were used for directionality.

\textbf{Fixed-lag sensitivity.} As a sensitivity analysis, fixed lags (lag $= 1, 3, 5$) were
examined for a 200-case subsample; GC directionality was consistent (EEG-preceding in 68--73\%
of cases across lags), supporting that the minimum-$p$-value lag selection does not materially
bias the directional conclusion.

% ── S4 ───────────────────────────────────────────────────────────────────────
\section{S4. Wavelet Implementation}

\texttt{pywt.cwt}; wavelet: \texttt{cmor1.5-1.0}; scales:
\texttt{np.arange(2, min(n/4, 150), 2)}; $dt = 10$~s. Smoothing:
\texttt{scipy.ndimage.uniform\_filter([5, 11])}.

% ── S5 ───────────────────────────────────────────────────────────────────────
\section{S5. Wavelet Edge-Effect Sensitivity}

Cone-of-influence (COI) edge effects affect $\approx$5\% of samples at record boundaries.
WC$_\text{overall}$ was recomputed after excluding the boundary 10\% of windows in each
record. Mean WC$_\text{overall}$ changed by $<$0.003 (from 0.799 to 0.801), confirming that
retained boundary data do not materially affect the results.

% ── S6 ───────────────────────────────────────────────────────────────────────
\section{S6. ADR--MAP Feature Pair Justification}

ADR--MAP was selected a priori. ADR correlates with BIS-defined anaesthetic depth (Pearson
$r \approx 0.7$ in prior studies), is continuous, and showed lower outlier rates than SEF95 or
spectral entropy in exploratory checks. MAP was selected as the primary haemodynamic variable
(the clinical monitoring target). For transparency, GC directionality for alternative pairs was
examined in a 200-case subsample: SEF95--MAP, 68\% EEG-preceding; spectral entropy--MAP, 65\%
EEG-preceding --- consistent with the primary ADR--MAP result (71\%).

% ── SUPPLEMENTARY FIGURES ────────────────────────────────────────────────────
\clearpage

\begin{figure}[H]
\centering
\includegraphics[width=0.72\linewidth]{supp_fig1_granger.png}
\caption{Mean Granger-causality F-statistics for EEG and haemodynamic feature pairs across
2,070 cases.}
\label{sfig:granger}
\end{figure}

\begin{figure}[H]
\centering
\includegraphics[width=\linewidth]{supp_fig2_quartiles.png}
\caption{Coupling component metrics by NCCI quartile. The component metrics were used to
construct NCCI; this ordering is a construction artefact and does not constitute independent
validation.}
\label{sfig:quartiles}
\end{figure}

\begin{figure}[H]
\centering
\includegraphics[width=0.75\linewidth]{supp_fig3_adr_response.png}
\caption{Case-level Alpha/Delta Ratio (ADR) response at surrogate haemodynamic excursions.
Provided for reference; the primary event analysis uses WCS (main-text Figure~4).}
\label{sfig:adr}
\end{figure}

\end{document}
